% Fragmento integrable. Preambulo anfitrion: amsmath, amssymb, longtable, hyperref.
\section{Analytic identification of the original centers}
\label{hmtfg:fragmento:identificacion-analitica-de-los-centros-originales}

\subsection{Degree-two realization and metric identification of the centers}
\label{hmtfg:clasificacion:8}

The equality between the first parameter of the infinite itinerary and the stable crossing, proved by first exit, provides
\begin{equation}
\boxed{a_*:=\min\Lambda_\tau=\lambda_*.}
\label{hmtfg:clasificacion:eq:H}
\end{equation}

The family, dodecaphase profile, returns and selector remain exactly as in sections~\ref{hmtfg:clasificacion:1}--\ref{hmtfg:clasificacion:5}. The external results used below recognize the dynamics of that already constructed object. Their universal values do not enter the selection of coefficients or parameters of the HMT family.

\textbf{Metric transfer theorem.} The identification \eqref{hmtfg:clasificacion:eq:H}, together with the previously proved entry, admissibility and transversality bounds, implies that the original sequence of first new roots satisfies, for certain constants \(C>0\) and \(0<\theta<1\),
\begin{equation}
\boxed{\lambda_*-\lambda_n=C\delta_{\mathrm F}^{-n}
\bigl(1+O(\theta^n)\bigr),}
\label{hmtfg:clasificacion:eq:14}
\end{equation}
where \(\delta_{\mathrm F}>1\) is the unstable eigenvalue of the doubling operator at the recognized fixed point. In particular,
\begin{equation}
\frac{\lambda_{n-1}-\lambda_{n-2}}
     {\lambda_n-\lambda_{n-1}}
=\delta_{\mathrm F}+O(\theta^n).
\label{hmtfg:clasificacion:eq:15}
\end{equation}
The sequence eventually coincides with the centers of the local continuation, once both indices refer to the same physical period. The proof retains the first-root selector of section~\ref{hmtfg:clasificacion:4}.

\subsubsection{Obtaining a common quadratic-like domain}
\label{hmtfg:clasificacion:8.1}

Write \(g\) for the renormalization fixed point and
\(\Gamma(\lambda)=R^4f_\lambda\). The analytic chart used in the certificate is
\[
f(x)=1-x^2V\!\left(\frac{x^2-1}{2.5}\right),\qquad
V(z)=\frac{u}{10}+\sum_{k\ge1}\nu_kz^k,
\]
with difference norm \(|\Delta u|+\|\Delta\nu\|_1\). The admissibility inequalities prove that \(g\) is even, analytic, unimodal, quadratic and of doubling type; moreover, \(g(x)=\varphi(x^2)\) with \(\varphi'<0\) in \([0,1]\).

Theorem 2.1 of de Faria–de Melo–Pinto, applied to the combinatorial set consisting solely of doubling, provides a one-point limit set with a quadratic-like extension. Since \(Rg=g\), the real convergence in that theorem identifies this point with \(g\). Lemma 3.2 provides, in an analytic neighborhood of \(g\), a finite renormalization iterate with a quadratic-like domain and uniformly positive modulus. The locations are sections 2.1–2.3, pp. 736–738, and section 3.1, Lemma 3.2, p. 745, of \href{https://annals.math.princeton.edu/wp-content/uploads/annals-v164-n3-p01.pdf}{de Faria–de Melo–Pinto, \emph{Global hyperbolicity of renormalization for \(C^r\) unimodal mappings}}.

To check membership in the neighborhood of that lemma, let \(\Omega_\epsilon\) be a sufficiently small complex neighborhood of \([-1,1]\), with
\[
\rho_\epsilon:=\sup_{x\in\Omega_\epsilon}
\left|\frac{x^2-1}{2.5}\right|<1,
\qquad M_\epsilon:=\sup_{x\in\Omega_\epsilon}|x|^2<\infty.
\]
The restriction of the certified chart satisfies
\begin{equation}
\|\Delta f\|_{\Omega_\epsilon}
\le M_\epsilon\left(\frac{|\Delta u|}{10}
          +\rho_\epsilon\|\Delta\nu\|_1\right)
\le C_\epsilon\|\Delta f\|_{\mathcal A}.
\label{hmtfg:clasificacion:eq:16}
\end{equation}
The certificate gives
\[
\|R^m\Gamma(\lambda_*)-g\|_{\mathcal A}
<0.001666(0.83996)^m.
\]
A fixed integer \(m\) is therefore chosen to place this image in the neighborhood of Lemma 3.2. Continuity of the finite composition provides an interval \(I_*\) around \(\lambda_*\) where the same operations are defined. If \(N\) is the iterate supplied by that lemma, we set
\begin{equation}
d=4+m+N,\qquad
\mathcal F_\lambda=R^df_\lambda,\quad \lambda\in I_*.
\label{hmtfg:clasificacion:eq:17}
\end{equation}
The family \(\mathcal F\) has quadratic-like representatives of degree two, on a common domain and with uniformly positive modulus. Its real-analytic dependence admits a local complexification of the parameter. The number \(d\) is fixed and counts the doublings performed before applying the parameter theorem.

This construction fixes a common domain and a finite number of returns \(d\), sufficient to apply the parameter theorem.

\subsubsection{Transport of certified transversality}
\label{hmtfg:clasificacion:8.2}

The direction of \(\Gamma\) satisfies the cone condition
\[
\|\dot\nu\|_1\le\tfrac14|\dot u|,\qquad \dot u\ne0.
\]
The certified differential bounds imply invariance of this cone and
\begin{equation}
|\dot u_{k+1}|\ge4.381|\dot u_k|
\label{hmtfg:clasificacion:eq:18}
\end{equation}
along the orbit of \(\Gamma(\lambda_*)\). All its points remain in the certified domain. In particular, the direction transported through the finite steps of \eqref{hmtfg:clasificacion:eq:17} retains a nonzero expanding component.

The hybrid class of the fixed point coincides with its stable manifold in the space of quadratic-like germs: \href{https://www.maths.tcd.ie/EMIS/journals/Annals/149_2/lyubich.pdf}{Lyubich, \emph{Feigenbaum–Coullet–Tresser universality and Milnor's hairiness conjecture}, Theorem 6.1, pp. 371–372}. Transversality remains valid on passing to that space, as can be checked without identifying different Banach norms. Indeed, if \(\partial_\lambda\mathcal F_{\lambda_*}\) were tangent to the hybrid class, it would be the tangent of an analytic disk contained in it. Uniform exponential convergence of renormalization on a compact subdisk, followed by Cauchy’s estimate in its parameter, would make the evaluation at \(x=1\) of its renormalized tangents converge to zero. However, in the preceding chart,
\[
\dot f(1)=-\dot u/10,
\]
and \eqref{hmtfg:clasificacion:eq:18} makes its absolute value grow. Both derivatives represent the same germ and have the same real evaluation. The contradiction proves
\begin{equation}
\partial_\lambda\mathcal F_{\lambda_*}
\notin T_{\mathcal F_{\lambda_*}}\mathcal H_{c_\infty}.
\label{hmtfg:clasificacion:eq:19}
\end{equation}
Thus, \(S=\{\mathcal F_\lambda\}\) is a complex-analytic transversal to the hybrid class of the doubling limit.

\subsubsection{Eventual identification of the first roots by their period}
\label{hmtfg:clasificacion:8.3}

Denote by \(c_r\) the canonical quadratic center of period \(2^r\), obtained through \(r\) successive doublings of the period-one center. This symbol denotes the parameter of the quadratic recognition family; it is distinct from the orbital functions \(c_j(\lambda)=f_\lambda^j(0)\) of the preceding sections. We write \(c_r^{\rm quad}\) to preserve the distinction.

The global continuation theorem and \eqref{hmtfg:clasificacion:eq:H} give \(\lambda_n\to\lambda_*\), hence \(\lambda_n\in I_*\) for sufficiently large \(n\). The classification and restrictive intervals of section~\ref{hmtfg:clasificacion:2} show that, when \(n>d\),
\begin{equation}
\mathcal F_{\lambda_n}=R^df_{\lambda_n}
\quad\text{has exact critical period }2^{n-d}
\quad\text{and itinerary }W_{n-d}0.
\label{hmtfg:clasificacion:eq:20}
\end{equation}
Its critical orbit remains in the invariant real interval, contained in the quadratic-like domain; its filled Julia set is therefore connected. Straightening preserves the postcritical orbit and the real branches. Following its restrictive returns gives \(n-d\) canonical doublings and, finally, a fixed critical point. The latter straightens to the center \(0\); the successive inverses of the straightening maps of the doubling copies give the unique center \(c_{n-d}^{\rm quad}\). This is the concrete application of the combinatorial parameterization and the straightening homeomorphism of the copies described in Lyubich, sections 5.1–5.3, pp. 362–364. We conclude
\begin{equation}
\mathcal F_{\lambda_n}\in\mathcal H_{c_{n-d}^{\rm quad}}.
\label{hmtfg:clasificacion:eq:21}
\end{equation}

Lyubich’s Theorem 7.4, pp. 387–388, gives a unique local intersection of \(S\) with each class \(\mathcal H_{c_r^{\rm quad}}\) for sufficiently large \(r\). Its assumptions of a priori complex bounds hold for the real quadratic doubling cascade, according to Theorem 5.6, p. 367, of the same work. Write \(\zeta_r\) for the parameter of that intersection. Equations \eqref{hmtfg:clasificacion:eq:20}–\eqref{hmtfg:clasificacion:eq:21}, convergence to the crossing and local uniqueness imply
\begin{equation}
\boxed{\lambda_n=\zeta_{n-d}\quad\text{for every }n\text{ sufficiently large}.}
\label{hmtfg:clasificacion:eq:22}
\end{equation}
The uniqueness used here concerns \textbf{each canonical hybrid class} in a fixed neighborhood of the crossing. It is not inferred merely from the existence of a local cascade or from uniqueness of the crossing with the stable sheet. Equation \eqref{hmtfg:clasificacion:eq:22} identifies the original selector, already defined globally, with the local intersections; it does not introduce a replacement selector.

In particular, let \(\mu_j\) be the local cascade of section \ref{hmtfg:escalado:8}, and let \(s\) be the fixed integer determined by its exact physical period:
\[
\operatorname{per}_{\rm crit}(f_{\mu_j})=2^{j+s}.
\]
In the presentation with curve \(R^4f_\lambda\) and period-two target, \(s=5\); the additional fixed steps used to localize the target modify \(s\). By \eqref{hmtfg:clasificacion:eq:21} and the same uniqueness,
\begin{equation}
\mu_j=\zeta_{j+s-d},\qquad
\boxed{\lambda_n=\mu_{n-s}\quad(n\text{ sufficiently large}).}
\label{hmtfg:clasificacion:eq:23}
\end{equation}
The shift \(d\) in the quadratic-like preparation and the shift \(s\) in the target periods are distinct objects. Identification is made through the physical period, not by equating those two integers.

\subsubsection{Power-law remainder through transverse holonomy}
\label{hmtfg:clasificacion:8.4}

Lyubich’s Lemma 7.3, p. 384, with proof on pp. 384–387, provides \(C^{1+\beta}\)-conformal regularity, for some \(\beta>0\), of the holonomy between transversals to the hybrid class of the Feigenbaum point. Its definition on p. 384 concerns the connectedness loci on the transversals; these contain, in particular, the centers considered here. We may reduce the exponent so that \(0<\beta\le1\).

Let \(h\) be the holonomy from \(S\) to the unstable manifold, and let \(z\) be an analytic coordinate linearizing the one-dimensional restriction of renormalization:
\begin{equation}
z(g)=0,\qquad z(Rq)=\delta_{\mathrm F} z(q).
\label{hmtfg:clasificacion:eq:24}
\end{equation}
The existence of this coordinate follows by applying the linearization of a repelling holomorphic germ, with multiplier \(\delta_{\mathrm F}>1\), to the analytic unstable manifold. The expanding direction of \eqref{hmtfg:clasificacion:eq:18}, cone invariance and the stationary return identify that multiplier with the positive unstable eigenvalue used in the certificate.

If \(q_r\) is the center of the class \(\mathcal H_{c_r^{\rm quad}}\) lying on the unstable manifold, compatibility of straightening and renormalization gives \(Rq_r=q_{r-1}\). This identity is the one used in the proof of Theorem 7.4, p. 388. By \eqref{hmtfg:clasificacion:eq:24}, there exists \(b\ne0\), absorbing any fixed initial index, such that
\begin{equation}
z(q_r)=b\delta_{\mathrm F}^{-r}
\label{hmtfg:clasificacion:eq:25}
\end{equation}
for every sufficiently large \(r\).

The regularity of \(h\) and transversality \eqref{hmtfg:clasificacion:eq:19} give, on the connectedness locus near \(\lambda_*\),
\begin{equation}
H(\lambda):=z\bigl(h(\mathcal F_\lambda)\bigr)
=A(\lambda-\lambda_*)
+O\!\left(|\lambda-\lambda_*|^{1+\beta}\right),
\qquad A\ne0.
\label{hmtfg:clasificacion:eq:26}
\end{equation}
Since \(H(\zeta_r)=b\delta_{\mathrm F}^{-r}\), the lower estimate
\(|H(\lambda)|\ge |A|\,|\lambda-\lambda_*|/2\) in a sufficiently small neighborhood first implies
\(|\zeta_r-\lambda_*|=O(\delta_{\mathrm F}^{-r})\). Substituting this bound into the remainder of \eqref{hmtfg:clasificacion:eq:26} yields
\begin{equation}
\zeta_r-\lambda_*=\frac bA\delta_{\mathrm F}^{-r}
+O\!\left(\delta_{\mathrm F}^{-(1+\beta)r}\right).
\label{hmtfg:clasificacion:eq:27}
\end{equation}
We apply \eqref{hmtfg:clasificacion:eq:22} with \(r=n-d\). The finite shift modifies the leading coefficient and the remainder constant. Since \(\lambda_n<\lambda_*\), we obtain \eqref{hmtfg:clasificacion:eq:14}, with
\[
C=-\frac bA\delta_{\mathrm F}^d>0,
\qquad \theta=\delta_{\mathrm F}^{-\beta}\in(0,1).
\]
Finally,
\[
\lambda_n-\lambda_{n-1}
=C(\delta_{\mathrm F}-1)\delta_{\mathrm F}^{-n}\bigl(1+O(\theta^n)\bigr),
\]
which proves \eqref{hmtfg:clasificacion:eq:15}. \(\square\)


\subsubsection{Quantitative scope of the theorem}
\label{hmtfg:clasificacion:8.5}

The identification of limits proved by first exit activates metric transfer for the global first roots of the uniform sector. The operator rate \(0.83996\) controls the renormalized maps on the stable sheet. The parameter rate \(\theta=\delta_{\mathrm F}^{-\beta}\) follows from the regularity of transverse holonomy. Both estimates retain their variable and domain.

The theorem establishes the existence of \(\beta,\theta,C\), the neighborhood \(I_*\), the finite number of returns \(d\) and an index \(n_0\) beyond which the centers coincide. Their individual numerical values are not assigned in this asymptotic statement. This formulation completes the limit and its power-law remainder; numerical evaluation of those witnesses is a distinct quantitative specialization.
